\documentclass[10pt,a4paper]{article}

\usepackage[margin=20mm]{geometry}
\usepackage{amsmath,amssymb,amsthm}

\theoremstyle{definition}
\newtheorem*{definition*}{Definition}
\newtheorem*{problem*}{Problem}

\newcommand{\CC}{\mathsf{CC}}
\newcommand{\CL}{\mathsf{CL}}
\newcommand{\NL}{\mathsf{NL}}
\newcommand{\NC}{\mathsf{NC}}
\newcommand{\TC}{\mathsf{TC}}
\newcommand{\SAC}{\mathsf{SAC}}
\newcommand{\Pclass}{\mathsf{P}}
\newcommand{\ZPP}{\mathsf{ZPP}}
\newcommand{\CSPACE}{\mathsf{CSPACE}}
\newcommand{\AC}{\mathsf{AC}}

\setlength{\parindent}{0pt}
\setlength{\parskip}{5pt}

\begin{document}

\begin{center}
  {\Large\bfseries Comparator Circuits versus Catalytic Logspace}
\end{center}

\section{Definitions}

\begin{definition*}[Comparator circuits and $\CC$]
A \emph{comparator gate} is the two-input, two-output Boolean function
\begin{equation}
  C:\{0,1\}^{2}\longrightarrow\{0,1\}^{2},
  \qquad
  C(p,q)=(p\wedge q,p\vee q).
\end{equation}
A comparator circuit on $m$ wires is specified by a finite sequence
\begin{equation}
  (i_1,j_1),\ldots,(i_t,j_t),
\end{equation}
where gate $(i_k,j_k)$ applies $C$ to the current contents of wires $i_k$ and
$j_k$. Each gate replaces those values by their minimum and maximum; it does
not create a copy of either input. Thus the model has fan-out one.

An \emph{annotated comparator circuit} additionally designates one output wire
and labels every input wire by a constant $0$ or $1$, an input variable $x_i$,
or a negated input variable $\neg x_i$. A family $\{C_n\}_{n\geq1}$ is
$\AC^0$-uniform when the description of $C_n$ is generated by a function in
uniform $\mathsf{FAC}^0$. The class $\CC$ consists of the relations computed by
$\AC^0$-uniform families of polynomial-size annotated comparator circuits.
This is equivalent to the definition as the languages logspace-reducible to
the comparator circuit value problem $\mathsf{CCV}$.
\end{definition*}

\begin{definition*}[Catalytic logspace and $\CL$]
A \emph{catalytic Turing machine} is a deterministic Turing machine equipped
with a read-only input tape, an ordinary read-write work tape, and an
additional read-write auxiliary tape. For an input $x$ and an arbitrary
string $a$ initially written on the auxiliary tape, let $M(x,a)$ denote the
resulting computation. The machine $M$ decides a language $L$ if, for every
$x$ and every $a$,
\begin{enumerate}
  \item $M(x,a)$ halts with the auxiliary tape restored exactly to $a$; and
  \item $M(x,a)$ accepts if and only if $x\in L$.
\end{enumerate}

For functions $S,S_a:\mathbb N\to\mathbb N$, let
$\CSPACE(S(n),S_a(n))$ be the class of languages decided with $O(S(n))$ cells
of ordinary work tape and $O(S_a(n))$ cells of auxiliary tape. Write
\begin{equation}
  \CSPACE(S(n))
  :=
  \CSPACE\bigl(S(n),2^{O(S(n))}\bigr)
\end{equation}
and define
\begin{equation}
  \CL:=\CSPACE(\log n).
\end{equation}
Equivalently, a $\CL$ machine has $O(\log n)$ ordinary work space and
$n^{O(1)}$ auxiliary space. No polynomial running-time bound is imposed.
\end{definition*}

\section{Best Known Result}

The unconditional containments relevant to the comparison are
\begin{equation}
  \NL\subseteq\CC\subseteq\Pclass,
  \qquad
  \NL\subseteq\TC^1\subseteq\CL\subseteq\ZPP.
\end{equation}
Thus $\NL$ is a nontrivial common subclass. Both classes also have the common
upper bound $\ZPP$, since
\begin{equation}
  \CC\subseteq\Pclass\subseteq\ZPP,
  \qquad
  \CL\subseteq\ZPP.
\end{equation}
These bounds do not imply a containment in either direction. In particular,
the $\CC$-complete problem $\mathsf{CCV}$ is not known to lie in $\CL$, and no
general simulation of catalytic machines by polynomial-size comparator
circuits is known.

Relativized $\NC$ and $\CC$ are incomparable in the length-preserving
functional-oracle model: for suitable $\alpha$ and $\beta$, not necessarily
the same oracle,
\begin{equation}
  \CC(\alpha)\not\subseteq\NC(\alpha),
  \qquad
  \NC^3(\beta)\not\subseteq\CC(\beta).
\end{equation}
This remains true for strictly $1$-Lipschitz oracles and supports, without
proving, the conjectured unrelativized incomparability.

The strongest circuit-class containment in $\CL$ stated in the source is
\begin{equation}
  \TC^1\subseteq\CL.
\end{equation}
Register-program simulations go beyond logarithmic depth, but use more than
logarithmic clean space. In particular,
\begin{equation}
  \NC^2
  \subseteq
  \CSPACE\left(
    O\left(\frac{\log^2n}{\log\log n}\right),
    n^{O(1)}
  \right),
\end{equation}
and, for every $\varepsilon>0$,
\begin{equation}
  \SAC^2
  \subseteq
  \CSPACE\left(
    O\left(\frac{\log^2n}{\log\log n}\right),
    2^{O(\log^{1+\varepsilon}n)}
  \right).
\end{equation}
Because the ordinary work-space bound in these simulations is
$\omega(\log n)$, they do not establish $\NC^2\subseteq\CL$ and do not resolve
the comparison with $\CC$. The relationship between $\CC$ and $\CL$ therefore
remains open.

\section{Problem Formalization}

For the uniform, unrelativized classes defined above, determine the
relationship between comparator circuits and catalytic logspace. Establish at
least one of the following alternatives:
\begin{enumerate}
  \item prove $\CC\subseteq\CL$;
  \item prove $\CL\subseteq\CC$; or
  \item prove that the classes are incomparable,
  \begin{equation}
    \CC\not\subseteq\CL
    \qquad\text{and}\qquad
    \CL\not\subseteq\CC.
  \end{equation}
\end{enumerate}


\end{document}
